\subsection{Application: examples of Noetherian rings}

LEMMA 1. Let A be a graded ring of type Z, whose graduation ( \(A_{n}\) ) is such that \(A_{n} = 0\) for all n \textless{} 0 or \(A_{n} = 0\) for all n \textgreater{} 0. Let M be a graded A-module of type Z. For M to be a Noetherian A-module, it is necessary and sufficient that every graded submodule of M be finitely generated.

As \(n \mapsto -n\) is an automorphism of the group \(\mathbf{Z}\) , we may restrict our attention to the case \(A_n = 0\) for all \(n > 0\) . Let \(A'\) and \(M'\) denote the ring \(A\) and the module \(M\) with the filtrations associated with their respective graduations (no. 1, Example 1), which are exhaustive and separated; the hypothesis on \(A\) implies that \(A'\) is discrete and hence complete. If \(E\) is a sub-A-module of \(M\) , the filtered \(A'\) -module \(E'\) obtained by giving \(E\) the induced filtration is Hausdorff and its filtration is exhaustive; moreover \(\text{gr}(E')\) is identified with a graded sub-A-module of \(M = \text{gr}(M')\) and hence is finitely generated by hypothesis. The conclusion then follows from Corollary 1 to Proposition 12 of no. 9.

THEOREM 2. Let A be a graded ring of type N, M a graded A-module of type N and \((\mathbf{A}_{n})\) and \((\mathbf{M}_{n})\) their respective graduations. Suppose that there exists an element \(a \in A_{1}\) such that \(A_{n} = A_{0}a^{n}\) and \(M_{n} = a^{n}M_{0}\) for all n \textgreater{} 0. Then, if \(M_{0}\) is a Noetherian \(A_{0}\) -module, M is a Noetherian A-module.

By virtue of Lemma 1 it is sufficient to prove that every graded submodule N of M is finitely generated. For all \(r \geqslant 0\) , let \(\mathbf{N}_r = \mathbf{N} \cap \mathbf{M}_r\) and let \(\mathbf{L}_r\) be the set of \(m \in \mathbf{M}_0\) such that \(a^r m \in \mathbf{N}_r\) . As

\[
a ^ {r} \mathrm{A} _ {0} \subset \mathrm{A} _ {r} = \mathrm{A} _ {0} a ^ {r}, \quad a ^ {r} \mathrm{A} _ {0} \mathrm{L} _ {r} \subset \mathrm{A} _ {0} a ^ {r} \mathrm{L} _ {r} \subset \mathrm{A} _ {0} \mathrm{N} _ {r} \subset \mathrm{N} _ {r}
\]

and hence the \(L_{r}\) are sub- \(A_{0}\) -modules of \(M_{0}\) ; moreover,

\[
a \mathrm{N} _ {r} \subset \mathrm{N} \cap a \mathrm{M} _ {r} = \mathrm{N} \cap \mathrm{M} _ {r + 1} = \mathrm{N} _ {r + 1}
\]

and hence the sequence \((\mathbf{L}_{r})_{r\geqslant0}\) is increasing. The hypothesis implies that there exists an integer \(n\geqslant0\) such that \(L_{r}=L_{n}\) for \(r\geqslant n\) . For each \(r\leqslant n\) , let \((m_{r,s})_{1\leqslant s\leqslant k_{r}}\) be a system of generators of the \(A_{0}\) -module \(L_{r}\) . We shall show that the elements \(a^{r}m_{r,s}\) for \(1\leqslant s\leqslant k_{r}, 0\leqslant r\leqslant n\) form a system of generators of the A-module N. As \(M_{r}=a^{r}M_{0}\) for all r, \(N_{r}=a^{r}L_{r}\) for all r by definition of \(L_{r}\) . Then, for \(r\leqslant n\) ,

\[
\mathbf {N} _ {r} = a ^ {r} \mathbf {L} _ {r} = \sum_ {s = 1} ^ {k _ {r}} a ^ {r} \mathbf {A} _ {0} m _ {r, s} \subset \sum_ {s = 1} ^ {k _ {r}} \mathbf {A} _ {0} a ^ {r} m _ {r s},
\]

and, for \(r > n\) ,

\[
\mathrm{N} _ {r} = a ^ {r} \mathrm{L} _ {n} = \sum_ {s = 1} ^ {k _ {n}} a ^ {r} \mathrm{A} _ {0} m _ {n, s} \subset \sum_ {s = 1} ^ {k _ {n}} \mathrm{A} _ {0} a ^ {r} m _ {n, s} \subset \sum_ {s = 1} ^ {k _ {n}} \mathrm{A} _ {0} a ^ {r - n}. (a ^ {n} m _ {n, s})
\]

which completes the proof (cf.~Exercise 10).

COROLLARY 1 (Hilbert's Theorem). For every commutative Noetherian ring C, the polynomial ring C{[}X{]} is Noetherian (cf.~Exercise 10).

COROLLARY 2. For every commutative Noetherian ring C and every integer n \textgreater{} 0, the polynomial ring C{[}X₁, \ldots, Xₙ{]} is Noetherian.

This follows from Corollary 1 by induction on n.

COROLLARY 3. If C is a commutative Noetherian ring, every finitely generated commutative C-algebra is a Noetherian ring.

Such an algebra is isomorphic to a quotient of a polynomial algebra \(\mathbf{C}[\mathbf{X}_1,\ldots ,\mathbf{X}_n]\) (§ 1, no. 1).

COROLLARY 4. Let A be a graded commutative ring of type N and let \((\mathbf{A}_{n})\) be its graduation. For A to be Noetherian, it is necessary and sufficient that \(A_{0}\) be Noetherian and that A be a finitely generated \(A_{0}\) -algebra.

The condition is sufficient by Corollary 3. Conversely, suppose A is Noetherian; \(m = \sum_{n \geqslant 1} A_n\) , which is an ideal of A, is then finitely generated; then A is a finitely generated \(A_0\) -algebra (§ 1, no. 2, Corollary to Proposition 1); on the other hand \(A_0\) , which is isomorphic to A/m, is a Noetherian ring.

COROLLARY 5. Let A be a commutative ring and m an ideal of A such that A/m is Noetherian, m/m² is a finitely generated (A/m)-module and A is Hausdorff and complete with the m-adic topology. Then gr(A) and A are Noetherian.

\(\operatorname{gr}(\mathbf{A})\) is an \((\mathbf{A} / \mathfrak{m})\) -module generated by \(\mathfrak{m} / \mathfrak{m}^2\) (no. 3, Example 3) and hence the ring gr(A) is Noetherian by Corollary 3. From this we deduce that A itself is Noetherian (no. 9, Corollary 2 to Proposition 12).

COROLLARY 6. For every commutative Noetherian ring C and every integer n \textgreater{} 0, the ring of formal power series C{[}{[}X₁, \ldots, Xₙ{]}{]} is Noetherian.

This follows from Corollary 5 and no. 6, Corollary to Proposition 6, for if m is the ideal of A = C{[}{[}X₁, \ldots, Xₙ{]}{]} consisting of the formal power series with no constant term, A/m is isomorphic to C and m/m² to the C-module Cⁿ.

\textbf{Remarks}\par

\begin{enumerate}
\def\labelenumi{(\arabic{enumi})}
\tightlist
\item
  Corollaries 2, 3 and 6 apply in particular if C is a commutative field.
\end{enumerate}

* (2) Let g be a Lie algebra over a commutative Noetherian ring C and suppose that g is a finitely generated C-module. Let the enveloping algebra U of g be given the increasing filtration \((\mathbf{U}_{n})\) defined in no. 3, Example 4. With the corresponding topology, U is discrete and hence Hausdorff and complete; the associated graded ring gr(U) is a finitely generated C-algebra, being a quotient of the symmetric algebra S(g), hence gr(U) is a Noetherian ring (Corollary 3) and we deduce that U is a left and right Noetherian ring (no. 9, Corollary 2 to Proposition 12). *

